\chapter{Hiện thực hệ thống}
\label{chap:hienthuc}

% TODO: toàn bộ chương này cần nhóm điền nội dung thực tế từ quá trình phát triển. Khung mục dưới đây bám theo Ch.3 để đảm bảo tính nhất quán thiết kế--hiện thực.

\section{Môi trường và công cụ phát triển}
\label{sec:moitruong}
% TODO: liệt kê ngôn ngữ, framework (frontend/backend), cơ sở dữ liệu, dịch vụ bên thứ ba (thanh toán, vận chuyển, lưu trữ đối tượng), công cụ AI/thư viện ML dùng cho các mô-đun tại mục~\ref{sec:thietkeAI}, công cụ quản lý mã nguồn/CI-CD nếu có.

\section{Kiến trúc triển khai}
\label{sec:trienkhai}
% TODO: mô tả hạ tầng triển khai thực tế (server/cloud, container hóa nếu có), đối chiếu với kiến trúc thiết kế tại mục~\ref{sec:kientruc}.

\section{Hiện thực các module chính}
\label{sec:hienthucmodule}

\subsection{Module xác thực \& quản lý tài khoản}
% TODO

\subsection{Module sản phẩm \& truy xuất nguồn gốc (QR)}
% TODO

\subsection{Module bán lẻ (B2C)}
% TODO

\subsection{Module đánh giá \& uy tín}
% TODO: hiện thực quy trình tại mục~\ref{subsec:danhgianghiepvu} (flowchart Hình~\ref{fig:flow-danh-gia}) -- thu thập đánh giá, kiểm duyệt, tính điểm uy tín tổng hợp.

\subsection{Module bán sỉ (B2B)}
% TODO

\subsection{Module quản lý tồn kho \& vận chuyển/khiếu nại}
% TODO: hiện thực quản lý tồn kho theo lô (FEFO) và tích hợp API với đơn vị vận chuyển 3PL -- mục~\ref{subsec:khobai}.

\subsection{Module quản trị}
% TODO

\subsection{Module AI}
\label{sec:hienthucAI}
% TODO: mô tả hiện thực cụ thể cho 4 mô-đun AI thiết kế tại mục~\ref{sec:thietkeAI} (gợi ý sản phẩm, dự báo nhu cầu/giá, phát hiện gian lận, trợ lý ảo) -- mô hình/thư viện sử dụng, dữ liệu huấn luyện, cách tích hợp vào hệ thống chính.

\section{Giải thuật và đoạn code quan trọng}
\label{sec:giaithuat}
% TODO: trích các đoạn code/giải thuật tiêu biểu (vd sinh mã QR, tính giá bậc theo MOQ, middleware RBAC, logic escrow) kèm giải thích ngắn gọn.
